\section{Transfer an agent design only with its causal question}
\label{exec18:reasoning-transfer}
An agent can give a convincing account of why a payment was recovered and
still choose an uneconomic intervention. The difficulty begins before model
selection: explaining a realised event, predicting a specified mechanism,
identifying an intervention distribution and choosing an action are different
requests. Two contrasting contributions clarify this distinction. A biomedical
position paper proposes combining language models, structured knowledge,
heterogeneous observations and formal causal tools. A philosophical study
instead supplies explicit neuron diagrams and asks a language model to identify
events and their causes. Both offer ideas for the banking product, but their
evidence supports different parts of its design.~\citep{exec18_reasoning_biomedical,exec18_reasoning_neuron}

\subsection{A domain architecture is a research agenda}
The selected initial version of Bazgir and colleagues' biomedical paper,
dated 22 May 2025, is explicitly a position paper. Its proposed workflow
combines literature retrieval, structured knowledge, heterogeneous observations
and external causal methods. It discusses control over consequential actions,
verification, task-specific evaluation, reproducibility and changing model
versions. These are useful architectural requirements. The paper's proposed
personalised and longitudinal applications are motivations, however; it does
not report a new matched experiment qualifying the complete causal agent or a
new identification argument for individual treatment decisions. Its claims
about constituent systems remain secondary reports that require those systems'
own technical originals.~\citep{exec18_reasoning_biomedical}

For banking, the useful transfer is therefore a division of responsibilities.
A language model can propose that a call-centre note contains a hardship
request, that an account event represents a treatment, or that a retrieved
document suggests a causal edge. Each proposal must retain the supporting
observation and its time. The scientific layer must separately decide whether
the proposed variable is measured consistently, whether the intervention is
well defined, and whether the available information identifies the requested
contrast. A graph of relationships extracted from documents does not by itself
supply a causal graph for the population or distinguish a predecision
covariate from a consequence of an earlier action. Calling an external
estimator can implement an approved analysis; it cannot make those scientific
premises true. This is the proposed banking interpretation of the architecture,
not a claim that the biomedical paper establishes the banking premises.

The source also illustrates why retrieval must preserve citation identity.
Its discussion of multimodal causal discovery attaches reference 7 to a
MATMCD passage, while its bibliography identifies reference 7 as Trialscope
and reference 17 as the multimodal-agent work. This internal mismatch does
not establish that the surrounding scientific statement is false, but it
prevents that reference number from serving as a verified technical path.
The product should preserve the passage and the mismatch, find the intended
original, and verify the claimed method there before treating it as available
support. An agent-written explanation of the mismatch would remain a
resolution proposal until the original is inspected.~\citep{exec18_reasoning_biomedical}

For a proposed servicing study, public or synthetic observations can test
whether the system preserves times, variables, graph restrictions and refusals
through this chain. Such a study can demonstrate the merit of building a
prototype without requiring bank deployment first. It cannot establish the
truth of a particular bank's undocumented assignment or audit process. The
proposal should therefore state what the public experiment would decide,
which scientific assumptions it stipulates, and what later institution-specific
evidence would change the decision. This boundary permits useful prototype
work while keeping the proposed transfer conditional.

\subsection{A supplied causal mechanism changes the evaluation task}
The selected Vervoort--Nikolaev manuscript has an internal dateline of
30 April 2025 and an arXiv submission of 17 June 2025. It presents textual
transcriptions of neuron diagrams, with stipulated activation, inhibition
and time-order rules. Its two questions ask whether the target event occurs
and what causes that occurrence or nonoccurrence. The model receives the
mechanism; it is not asked to identify a causal graph from a sampled
observational distribution. Performance on this task therefore tests a
particular form of reasoning about supplied mechanisms. It cannot be read as
performance on statistical discovery or intervention identification from
unrestricted data.~\citep{exec18_reasoning_neuron}

The proposed labeling rule, DEF-1, begins with a counterfactual comparison
but introduces additional treatment of bifurcating paths. In selected cases it
retains factual blocking on off-path routes even when the candidate cause is
changed. The authors explicitly describe this as not fully implementing the
ordinary removal of that event throughout the mechanism. They also limit
the rule to a class of diagrams, discuss an ambiguous case requiring a
possible extension, and leave its relation to other causal theories open.
For the banking product this is a candidate attribution convention with a
declared domain, not a universal replacement for an intervention operator.
Implementing a library operation under this convention would require its
own formal specification and conformance tests. The example below does not
implement or validate DEF-1.~\citep{exec18_reasoning_neuron}

The evaluation evidence requires equally careful interpretation. The paper
describes the experiments as illustrations, explicitly distinguishes them
from large statistical studies, and reports that the displayed early tables
came from a single run without feedback. Its 25 selected diagrams also
helped develop the proposed definition. The prose reports 13 of 25 fully
correct answers and two partially correct answers for the selected GPT-4
exercise; these are historical, author-reported counts under that rubric,
not an estimate of general causal competence. The later model comparison
changes the prompt to request short answers and initial causes. Its result
cannot isolate a model-change effect from the changed evaluation contract.
No current model ranking follows from these experiments.~\citep{exec18_reasoning_neuron}

Even a small benchmark needs a faithful transcript and an explicit grading
unit. In the first diagram's recorded answer, the explanation assigns the
blocking of the backup path to A although the prompt assigns it to C; the
answer is nevertheless marked correct. The diagram 17 transcription asks
whether G occurs and then asks for the causes of E. These are local
transcript and target inconsistencies, not a basis for recomputing all scores
from extracted text. They show why an exact-output evaluator must distinguish
a correct final event, a correct cause list and a faithful explanation. A
reuse decision should preserve the originals, record disputed labels and
adjudicate the target before admitting the case. Diagram images, underlined
table annotations and full score reproduction need their own inspection.~\citep{exec18_reasoning_neuron}

\subsection{A servicing route with backup: attribution versus value}
\label{exec18:reasoning-example}
Consider a deliberately stipulated servicing mechanism. Let $A\in\{0,1\}$
denote availability of a backup contact, and $C\in\{0,1\}$ activation of a
priority service. At the next stage, $B$ records whether the backup actually
runs and $D$ whether the priority route delivers. Recovery $Y$ follows the
Boolean equations
\begin{equation}
 B=A(1-C),\qquad D=C,\qquad Y=\max\{B,D\}.
 \label{exec18:reasoning-mechanism}
\end{equation}
These are authored assumptions, not estimated account behavior. They encode
perfect success of an executed route and suppression of the backup while
priority is active. In the factual state $A=C=1$, the realised path is
$D=1$, $B=0$, $Y=1$. A faithful narrative can identify the executed priority
route. Under an intervention setting only $C=0$, however, the structural
equation for $B$ is still active: $B=1$ and $Y=1$. Removing priority alone
does not remove recovery in this state.~\citep{exec18_reasoning_reference}

If a different comparison also forces the backup to remain off, the outcome
is instead $Y=0$. With the subscript describing the intervention,
\begin{equation}
 Y_{do(C=0)}\bigm|_{A=1}=1,\qquad
 Y_{do(C=0,B=0)}\bigm|_{A=1}=0.
 \label{exec18:reasoning-interventions}
\end{equation}
The second expression changes two mechanisms. It may be relevant to a
declared attribution question, but it is not the outcome of disabling only
the priority service. A result object must retain which comparison was
requested. Replacing an intervention by a compound action that holds a
downstream route fixed silently changes the scientific question. Conversely,
a zero effect of disabling one route does not erase the fact that the route
was executed. The explanatory and intervention answers can differ without
logical contradiction.

To make the distinction economic, suppose a recovery is worth USD~200,
priority costs USD~40 when activated, and an executed backup costs USD~10.
The stipulated net value per case is
\begin{equation}
 N=200Y-40C-10B\quad\text{USD}.
 \label{exec18:reasoning-money}
\end{equation}
At $A=1$, priority produces USD~160 net and the backup-only route USD~190.
Priority's incremental net value is therefore minus USD~30 although it
provides the factual recovery path. When $A=0$, priority gives USD~160 while
no priority gives zero. If backup availability is not known before action
but its stipulated probability is $p$, averaging these two states gives
\begin{equation}
 \mathbb{E}[N_{do(C=1)}-N_{do(C=0)}]
 =160-190p\quad\text{USD per case}.
 \label{exec18:reasoning-threshold}
\end{equation}
Under these assumptions priority has positive incremental value precisely
when $p<16/19$, with indifference at equality. If availability is actually
observed before action, the appropriate conditional decision differs from
this population decision. The information clock must therefore accompany
the value calculation. Exact enumeration of all four $(A,C)$ states and
rational checks at $p=0,1/2,16/19,1$ reproduce the result. Neither the
mechanism nor $p$ nor the monetary inputs were learned from bank data; these
calculations establish an internal distinction, not bank ROI.~\citep{exec18_reasoning_reference}

\subsection{Admit the question and the rubric together}
A useful public evaluation can now separate tasks that a single accuracy
score would obscure. One task checks propagation of a supplied mechanism.
Another checks the stated attribution convention. A third asks for a
specific intervention and requires the correct action set. A fourth asks
for a value comparison under a stated predecision information set. An
answer that computes the compound intervention correctly but returns it for
the single-action query fails that contract. An answer that requests the
missing attribution convention may be scientifically preferable to an
unqualified list of causes.

The next prototype design should freeze the target, intervention semantics,
grading rubric and explanation requirements before evaluation. New graph
families should be held apart from those used to design the rules; word
changes alone do not make a familiar mechanism a new causal case. A public
corpus can retain disputed examples with an adjudication flag, rather than
silently treating one philosophical judgment as uncontested truth. Matched
comparisons should retain the same target, prompt conditions, tool access
and output budget, while reporting invalid, unsupported and incomplete
outputs separately. This is a proposed route to credible prototype evidence.
It leaves open the broader coverage of the literature, executable package
admission, live performance and the validity of institutional causal premises.
